\documentclass[11pt,a4paper]{article}
\usepackage[T1]{fontenc}
\usepackage[utf8]{inputenc}
\usepackage[french]{babel}
\usepackage{lmodern}
\usepackage{geometry}
\usepackage{amsmath,amssymb}
\usepackage{hyperref}
\geometry{margin=25mm}
\title{Optimisation d'une métasurface par FEM élastique transitoire\\
\large Réseau de géométrie et adjoint discret --- V4}
\author{Note technique}
\date{\today}
\begin{document}
\maketitle
\tableofcontents

\section{Objet et géométrie}
On recherche une carte de plots sur les deux faces d'une plaque pour
concentrer une onde supposée incidente depuis la droite sur une ligne
\(x=x_f\). La plaque et les plots sont décrits par des éléments
hexaédriques élastiques. Le réseau paramètre la géométrie, pas le champ :
ce n'est donc pas un PINN classique. La source ponctuelle représente une
onde retour supposée ; aucun défaut réel n'est simulé.

Dans la zone optimisable, le réseau produit une hauteur \(h_e\in[0,H]\)
en unités de couches, avec \(H=\texttt{MAX\_HEIGHT}\). Les plots sont des
colonnes pleines symétriques sur les faces supérieure et inférieure.
Une même colonne occupe toute la base d'un élément du maillage dans le plan.

\section{Réseau et périodes}
Deux paramètres appris \(a_x,a_y\) définissent des périodes nominales bornées
par une sigmoïde :
\[
p_x=p_x^{\min}+(p_x^{\max}-p_x^{\min})\sigma(a_x),\qquad
p_y=p_y^{\min}+(p_y^{\max}-p_y^{\min})\sigma(a_y).
\]
En chaque centre d'élément, le réseau reçoit les coordonnées normalisées
ainsi que les sinus et cosinus des phases périodiques dans les deux
directions. Il utilise deux couches cachées \(\tanh\), puis une sigmoïde.
Sa sortie est multipliée par deux enveloppes périodiques et par \(H\).
Les symétries demandées sont appliquées par moyenne avec les cartes
miroirs. Ces périodes sont nominales : l'arrondi et la modulation spatiale
peuvent modifier les espacements effectifs.

\section{Relaxation des couches}
Une hauteur continue permet de différencier le calcul. Pour une couche
d'indice \(k=0,\ldots,H-1\), on définit sa fraction de remplissage
\[
\phi_{ek}=\max(0,\min(1,h_e-k)),\qquad
\rho_{ek}=\varepsilon+(1-\varepsilon)\phi_{ek},
\]
où \(\varepsilon=\texttt{V4\_VOID\_FRACTION}\in(0,1)\).
La rigidité et la masse élémentaires des plots sont multipliées par
\(\rho_{ek}\) ; les deux faces utilisent la même fraction. La plaque
reste pleine. Les masses sont diagonalisées par sommation élémentaire.
Cette interpolation maintient un maillage et des masses strictement
positives, mais le vide résiduel et les couches intermédiaires ne sont pas
physiques et peuvent introduire des modes parasites. Aux points
non différentiables \(h_e=k\), l'implémentation prend une dérivée nulle.

Les matrices assemblées sont notées \(K(h)\) et \(M(h)\), avec \(M\)
diagonale. Les équations discrétisées sont
\[
 M(h)\ddot u(t)+K(h)u(t)=f(t).
\]
La source est une impulsion sinusoïdale fenêtrée sur le nombre de cycles
choisi. Les bords sont libres. Les efforts sur les faces peuvent favoriser
une parité A ou S sans garantir une excitation modale pure.

\section{Propagation et ligne focale}
Avec \(\Delta t\) fixe, le schéma explicite est
\[
u_{n+1}=2u_n-u_{n-1}
 +\Delta t^2M^{-1}(f_n-Ku_n),\qquad u_{-1}=u_0=0.
\]
Le pas commun aux configurations est borné de façon conservatrice à
partir des sommes absolues des lignes des matrices élémentaires divisées
par les masses nodales élémentaires. Il est donc utilisable aussi bien
pour la référence que pour la carte optimisée.

La ligne contient les nœuds \((x_f,y_j)\), de \(y=0\) à \texttt{WIDTH},
avec \(x_f\) arrondi au nœud le plus proche. Les poids sont
\[
w_j=\frac{q_j\exp[-(y_j-y_c)^2/(2\sigma_f^2)]}
{\sum_\ell q_\ell\exp[-(y_\ell-y_c)^2/(2\sigma_f^2)]},
\quad y_c=\frac{\texttt{WIDTH}}2,\quad
q_j=\begin{cases}\frac12&\text{aux extrémités},\\1&\text{ailleurs}.\end{cases}
\]
La largeur \(\sigma_f=\texttt{FOCAL\_SIGMA}\) vaut initialement 50 mm.
Le critère, sur les instants de \texttt{FOCUS\_WINDOW}, est
\[
G(h)=\frac{\displaystyle\sum_jw_j\sum_{n\in W}u_{z,j,n}(h)^2\Delta t}
{\displaystyle\sum_jw_j\sum_{n\in W}u_{z,j,n}(0)^2\Delta t},
\qquad J(h)=G(h)-\lambda\frac{\operatorname{moyenne}(h)}H.
\]
La référence \(h=0\) est une plaque nue sans matériau résiduel ; \(\lambda\)
est \texttt{V4\_HEIGHT\_PENALTY}. La grandeur intégrée est \(u_z^2\), non
l'énergie mécanique totale. Un rapport élevé ne garantit ni une intensité
uniforme le long de la ligne ni une tache étroite selon \(x\).

\section{Adjoint discret}
Posons \(A=2I-\Delta t^2M^{-1}K\),
\(B=-I\), et \(\ell_n(u_n)\) la contribution de l'instant \(n\) au
gain normalisé. La récurrence adjointe, avec
\(\lambda_N=\lambda_{N+1}=0\) après le dernier état, est
\[
\lambda_n=\nabla_{u_n}\ell_n+A^{\mathsf T}\lambda_{n+1}
                 +B^{\mathsf T}\lambda_{n+2}.
\]
Le gradient d'une hauteur, hors pénalité, s'écrit alors
\[
\frac{\partial G}{\partial h_e}
=-\Delta t^2\sum_n
\left(\frac{\lambda_{n+1}}{M}\right)^{\mathsf T}
\left[
\frac{\partial K}{\partial h_e}u_n+
\frac{\partial M}{\partial h_e}a_n
\right],
\quad
a_n=M^{-1}(f_n-Ku_n).
\]
La division par \(M\) est composante par composante. Les dérivées locales
de \(K\) et \(M\) viennent de la couche partiellement remplie sur chaque
face. Le gradient traverse ensuite le réseau et les périodes par
différentiation automatique ; Adam met à jour ses paramètres.

Les états complets ne sont pas tous stockés : des checkpoints réguliers
permettent de reconstruire chaque segment pendant le passage adjoint.
\texttt{V4\_CHECKPOINT\_INTERVAL} règle le compromis mémoire/recalcul.
Les calculs élastiques restent sur CPU, même si le réseau utilise un GPU.

\section{Sélection de la géométrie entière}
Tous les \texttt{V4\_DISCRETE\_EVERY} pas, les hauteurs sont arrondies
en couches entières. Une simulation élastique complète, sans matériau
résiduel, recalcule le critère avec le même pas de temps et la même
référence de plaque nue. La meilleure carte selon \(J\) est conservée.
Si aucune carte entière évaluée ne dépasse l'objectif de la plaque nue,
la plaque nue est conservée. La meilleure carte relaxée n'est donc pas
nécessairement la carte finale.

\section{Sorties et limites}
\texttt{design.npz} contient la carte entière et l'historique ;
\texttt{optimisation.json} rapporte son gain élastique et ses paramètres ;
\texttt{convergence.png} compare gains relaxés et entiers ;
\texttt{metasurface\_3d.png} montre les hauteurs sur les deux faces.
La vérification séparée produit les signaux, les profils de ligne et une
animation, avec notamment \texttt{verification\_3d.json} et
\texttt{ligne\_focale\_3d.png}.

La vérification utilise le même modèle physique que l'optimisation et ne
constitue pas une validation indépendante. L'optimisation est coûteuse ;
l'adjoint ne garantit pas un optimum global et la relaxation peut différer
fortement de la géométrie entière. Le maillage initial correspond à
environ quatre éléments par longueur d'onde A0 théorique selon \(x\) :
des études de convergence spatiale et temporelle, un autre solveur et une
validation expérimentale sont nécessaires avant toute conclusion physique.
\end{document}
